\documentclass[11pt]{article}
\usepackage[margin=1in]{geometry}
\usepackage{amsmath, amssymb, amsthm}
\usepackage{hyperref}

\newtheorem{theorem}{Theorem}
\newtheorem{definition}{Definition}
\newtheorem{proposition}{Proposition}

\title{Rubinstein Bargaining: Alternating Offers, Patience, and the Nash Solution}
\author{Abhi Wadhwa}
\date{}

\begin{document}
\maketitle

\begin{abstract}
Notes on Rubinstein's alternating-offers bargaining model and its connection to the Nash bargaining solution. I derive the unique subgame perfect equilibrium, explain why patient players get more, work through the finite-horizon case via backward induction, and show how the Rubinstein outcome converges to the Nash bargaining solution as offers become frequent. The main takeaway: \textbf{patience is power}, and the structure of sequential bargaining pins down a unique outcome, resolving the indeterminacy of the static Nash bargaining problem.
\end{abstract}

\section{The Problem}

Two players want to split a surplus of size 1 (think: a dollar, a pie, whatever). They take turns making offers. Player~1 proposes a split $(x, 1-x)$, and Player~2 either accepts (game over) or rejects. If Player~2 rejects, they make a counter-offer $(y, 1-y)$, and Player~1 accepts or rejects. This goes on, potentially forever.

The catch: both players are impatient. A dollar today is worth more than a dollar tomorrow. Player $i$ has discount factor $\delta_i \in (0, 1)$, meaning that receiving payoff $u$ in round $t$ is worth $\delta_i^t \cdot u$ in present value.

Before Rubinstein's 1982 paper \cite{rubinstein1982}, the standard approach to bargaining was Nash's axiomatic theory \cite{nash1950}, which tells you what the outcome ``should'' be based on fairness axioms but doesn't model the negotiation process. Rubinstein showed that you can actually model the negotiation as a game and get a \textbf{unique} prediction.

\section{The Infinite-Horizon SPE}

Here's where it gets good. Despite the game having infinitely many rounds and infinitely many possible offers, there's a \textbf{unique subgame perfect equilibrium}. The derivation is one of my favorite arguments in game theory.

The trick is to exploit the \textbf{stationarity} of the game. After any rejection, the game looks exactly like the original game (just with the roles swapped and payoffs discounted by one period). This means the continuation value after a rejection is just the equilibrium value of the ``same'' game with players swapped.

Let $x^*$ be Player~1's share in the SPE when Player~1 proposes, and let $y^*$ be Player~2's share when Player~2 proposes. In equilibrium, the responder is exactly \textbf{indifferent} between accepting and rejecting:

When Player~1 proposes and Player~2 responds:
\[
1 - x^* = \delta_2 \cdot y^*
\]
Player~2 gets $1 - x^*$ now if they accept, or waits one round and gets $y^*$ (discounted by $\delta_2$) if they reject and propose next.

When Player~2 proposes and Player~1 responds:
\[
1 - y^* = \delta_1 \cdot x^*
\]
Same logic, roles reversed.

Now solve this system. From the second equation, $y^* = 1 - \delta_1 x^*$. Plugging into the first:
\[
1 - x^* = \delta_2(1 - \delta_1 x^*) = \delta_2 - \delta_1 \delta_2 x^*.
\]
Rearranging:
\[
x^*(1 - \delta_1 \delta_2) = 1 - \delta_2,
\]
so
\[
\boxed{x^* = \frac{1 - \delta_2}{1 - \delta_1 \delta_2}}.
\]

Player 2's equilibrium share is:
\[
1 - x^* = \frac{\delta_2(1 - \delta_1)}{1 - \delta_1 \delta_2}.
\]

This is clean, and let me highlight a few things that I think are really interesting.

\textbf{Agreement is immediate.} In equilibrium, Player~1 makes an offer in round 1 and Player~2 accepts. There's never any delay, even though the game could go on forever. The ``threat'' of future bargaining is what pins down the price, but the threat is never executed.

\textbf{First-mover advantage.} Player~1 gets $x^*$, which is strictly more than $1/2$ whenever $\delta_2 < 1$ (and $\delta_1 \delta_2 < 1$, which is always true). There's a real advantage to going first, but it vanishes as the players become perfectly patient.

\textbf{Patient players get more.} This is the big insight: as $\delta_1$ increases (Player~1 becomes more patient), $x^*$ increases. As $\delta_2$ increases, $x^*$ decreases. Your bargaining power comes from your willingness to wait.

To see why patience matters intuitively: if you're patient, rejecting a bad offer is cheap for you (the future isn't much worse than the present). So your opponent has to offer you more to get you to accept now. Conversely, if you're impatient, you're desperate to settle quickly, so your opponent can low-ball you.

\section{Symmetric Case and Limits}

When $\delta_1 = \delta_2 = \delta$, the formula simplifies beautifully:
\[
x^* = \frac{1 - \delta}{1 - \delta^2} = \frac{1}{1 + \delta}.
\]

As $\delta \to 1$: $x^* \to 1/2$. \textbf{Patient, symmetric players split evenly.} The first-mover advantage vanishes because when $\delta$ is close to 1, waiting one round costs almost nothing, so neither side has leverage from the timing.

As $\delta \to 0$: $x^* \to 1$. Completely impatient players---the first mover takes everything because the responder values future payoffs at essentially zero.

\section{Finite Horizon}

With $T$ total rounds, we solve by backward induction. This is a nice exercise that connects to the infinite-horizon result.

\textbf{Round $T$} (last round, say Player~1 proposes if $T$ is odd): The proposer offers $(1, 0)$ and the responder accepts (anything non-negative). Proposer gets everything.

\textbf{Round $T-1$}: The proposer (Player~2 if $T$ is odd) must offer the other player at least their discounted continuation value. Player~1's continuation value from being the proposer next round is $\delta_1 \cdot 1 = \delta_1$. So Player~2 offers Player~1 exactly $\delta_1$ and keeps $1 - \delta_1$.

Continuing backward, you get a sequence of offers that alternates between the players, with each proposer giving the responder their discounted continuation value and keeping the rest.

The key observation is that as $T \to \infty$, the finite-horizon offers converge to the infinite-horizon SPE shares. The convergence is geometric, with rate $\delta_1 \delta_2$.

\section{The Nash Bargaining Solution}

Now let me bring in Nash's axiomatic approach from the other direction. The \textbf{Nash bargaining solution} (NBS) picks the split that maximizes the \textbf{Nash product}:
\[
\max_{(u_1, u_2)} \; (u_1 - d_1)^\alpha (u_2 - d_2)^{1-\alpha}
\]
subject to feasibility and individual rationality, where $(d_1, d_2)$ is the disagreement point (what each player gets if no deal is reached) and $\alpha \in (0,1)$ is Player~1's bargaining power.

For the simple case of splitting a surplus of size 1 with disagreement point $(0, 0)$, the NBS is:
\[
u_1^* = \alpha, \qquad u_2^* = 1 - \alpha.
\]

The NBS is characterized by four axioms:
\begin{enumerate}
\item \textbf{Pareto efficiency}: no feasible outcome makes both players better off.
\item \textbf{Symmetry}: symmetric problems get symmetric solutions (when $\alpha = 1/2$).
\item \textbf{Independence of irrelevant alternatives}: removing non-optimal feasible outcomes doesn't change the solution.
\item \textbf{Invariance to affine transformations}: rescaling utilities doesn't change the outcome (in utility terms).
\end{enumerate}

Nash (1950) showed these four axioms uniquely pin down the Nash product maximizer. This is beautiful but also kind of unsatisfying---it tells you what the answer is but not \emph{why} the players would arrive there through an actual negotiation process.

\section{The Rubinstein--Nash Convergence}

This is the punchline that ties everything together. Rubinstein's sequential game and Nash's axiomatic solution aren't competing theories---they're the same thing in the limit.

\begin{theorem}[Binmore, Rubinstein, and Wolinsky, 1986]
As $\delta_1, \delta_2 \to 1$ (offers become frequent / players become patient), the Rubinstein SPE converges to the asymmetric Nash bargaining solution with bargaining powers
\[
\alpha = \frac{\ln \delta_2}{\ln \delta_1 + \ln \delta_2}.
\]
\end{theorem}

When $\delta_1 = \delta_2$, this gives $\alpha = 1/2$---the symmetric NBS, which is the even split. When the discount factors differ, the more patient player (higher $\delta$) gets a larger bargaining power $\alpha$ (or $1-\alpha$, depending on which player they are).

After playing with this for a while, I think the convergence result is the deepest insight of the whole theory. It says that Nash's axioms aren't just abstract philosophy---they capture the limiting behavior of a concrete, fully specified negotiation game. \textbf{The ``bargaining power'' parameter $\alpha$, which Nash left unspecified, is determined by the relative patience of the players.} Patience is the microfoundation for bargaining power.

The technical argument goes roughly like this: write $\delta_i = e^{-r_i \Delta}$ where $r_i$ is player $i$'s discount rate and $\Delta$ is the time between offers. As $\Delta \to 0$ (offers become continuous), the SPE share becomes
\[
x^* = \frac{1 - e^{-r_2 \Delta}}{1 - e^{-(r_1 + r_2)\Delta}} \approx \frac{r_2 \Delta}{(r_1 + r_2)\Delta} = \frac{r_2}{r_1 + r_2},
\]
which is exactly the asymmetric NBS with $\alpha = r_2/(r_1 + r_2)$. The approximation is first-order in $\Delta$, so the convergence is tight.

\section{Wrapping Up}

The Rubinstein model is a cornerstone of bargaining theory because it resolves the indeterminacy problem. With Nash bargaining, you have axioms but no process. With Rubinstein, you have a process that delivers a unique outcome and connects back to the axioms in the limit. The model makes a sharp, testable prediction: \textbf{patient players get more}, agreement is immediate, and the split depends on the ratio of discount rates.

The practical implications extend well beyond the stylized two-player model. Labor negotiations, international trade deals, legislative bargaining---any situation where parties alternate proposals and are sensitive to delay can be analyzed through this lens.

\begin{thebibliography}{9}
\bibitem{rubinstein1982} A.~Rubinstein, ``Perfect equilibrium in a bargaining model,'' \textit{Econometrica}, vol.~50, no.~1, pp.~97--109, 1982.

\bibitem{nash1950} J.~F. Nash, ``The bargaining problem,'' \textit{Econometrica}, vol.~18, no.~2, pp.~155--162, 1950.

\bibitem{binmore1986} K.~Binmore, A.~Rubinstein, and A.~Wolinsky, ``The Nash bargaining solution in economic modelling,'' \textit{RAND Journal of Economics}, vol.~17, no.~2, pp.~176--188, 1986.

\bibitem{osborne1990} M.~J. Osborne and A.~Rubinstein, \textit{Bargaining and Markets}, Academic Press, 1990.

\bibitem{muthoo1999} A.~Muthoo, \textit{Bargaining Theory with Applications}, Cambridge University Press, 1999.
\end{thebibliography}

\end{document}
